\section{代码复用与 DRY}

在 Vibe Coding 场景下，模型受限于 Context Window，容易只看部分代码就开始修改，导致工具函数（Utils）出现大量冗余和重复代码。

Superpowers 在 \texttt{Writing Plans} 和 Code Review 中反复强调：

\begin{itemize}
    \item \textbf{DRY}（Don't Repeat Yourself）：杜绝重复代码。
    \item \textbf{YAGNI}（You Aren't Gonna Need It）：坚决杜绝过度设计和过度抽象。
\end{itemize}

此外，Superpowers 还提供了 \texttt{Writing Skills} 这一 Skill，强调将解决问题的过程提炼为\textbf{可复用的模式}，而非一次性脚本。

\subsection{本章小结}

DRY 与 YAGNI 是 Superpowers 贯穿始终的两大原则，前者防止冗余，后者防止过度工程化，两者构成代码质量的动态平衡。


